% Source-backed owned SDK purpose-key table; prior root modules are retained.
\subsection{An owned SDK purpose-key table}
\label{sec:sdk-purpose-connection}

The native SDK already derives several application keys from one combined
secret. Its \texttt{SharedSecret::derive\_key} operation takes a purpose,
protocol label and context, and returns a new \texttt{DerivedKey} owner.
The same source owner supports repeated calls. The pinned source exercises
all five purposes at both peers and checks their corresponding outputs and
lifecycle behavior. This supplies a concrete application of the finite
entry-table argument with the SDK's own information encoding.

\begin{lemma}[Information-length form of the entry comparison]
\label{lem:entry-information-length}
For the finite service $R_{r,I}=F(F(r,S),I\Vert\mathtt{01})$, let the
complete available information strings in class $i$ have length at most
$A_i$ bytes. Put $L_i=A_i+1$ and $\lambda_i=\ell(L_i)$, within the retained
SHA-256 padding domain. Use the descriptor classes and finite budgets of
Theorem~\ref{thm:multiple-root-entry}. With the parent contribution used
through the Extract interface and each PRK through its Expand row, the
whole interactive service is within output-bit gap
\begin{equation}
 \Delta_{\rm ent}=\min\!\left\{1,
 \varepsilon_{\rm Hjoint}(C,0,q_S,32)
 +\sum_{i=1}^{C}\varepsilon_{\rm Hrow}(N_i,q_i,L_i)\right\}
 \label{eq:entry-information-length}
\end{equation}
of the consistent independent table indexed by $(\Gamma(r),I)$.
The continuation may process each returned value through its declared
interface. Conservative complete ordinary-query budgets are
\[
 q_S=B+2D+\sum_i(2+\lambda_i)N_i,\qquad
 q_i=B+2D+\sum_{k>i}(2+\lambda_k)N_k.
\]
\end{lemma}
\begin{proof}
Use the column replacement and first-class-allocation row hybrids of
Theorem~\ref{thm:multiple-root-entry}. At row $i$, every queried message
has length at most $L_i$. The map $I\mapsto I\Vert\mathtt{01}$ is injective
also for variable-length strings. Hence the final private row gives exactly
the full-information table. Each actual Expand costs at most
$2+\ell(L_i)$ compression calls, yielding the displayed budgets.
The reductions receive the root-row replies and run the entire declared
continuation from them. This proves the table comparison for that
continuation. The native 97-byte case has $L_i=98$ and $\lambda_i=3$.
\end{proof}

\paragraph{The actual SDK descriptor.}
Let $\sigma=\texttt{QPeriapt-SDK-HKDF-SHA256-v1}$, a 27-byte string, and
$r_0=\sigma\Vert\mathtt{00}^5$, its 32-byte zero extension. HMAC pads both $\sigma$ and $r_0$ to the same
64-byte key block~\cite{rfc2104}, so $F(\sigma,S)=F(r_0,S)$ for every $S$ and every $h$.
The complete SDK information string is
\begin{equation}
\begin{split}
 I={}&D_{\rm sdk}\Vert\operatorname{BE32}(1)
 \Vert\operatorname{BE32}(2)\Vert\mathit{policyState}_{36}
 \Vert\mathit{rootDigest}_{32}\\
 &\Vert\operatorname{BE32}(p)
 \Vert\operatorname{BE16}(|a|)\Vert a
 \Vert\operatorname{BE32}(|c|)\Vert c
 \Vert\operatorname{BE16}(32),
\end{split}
 \label{eq:sdk-purpose-info}
\end{equation}
where $D_{\rm sdk}=\texttt{QPeriapt-SDK-Key-v1}$ has 19 bytes. The fixed
suite code $1$ denotes ML-KEM-768+X25519 and profile code $2$ denotes
ContextBound. Here $p$ is one
of the five purpose codes, $a$ is the protocol label, and $c$ is the context.
The policy state is its 4-byte version followed by its exact-document
32-byte digest; the root digest is SHA-256 of the runtime's pinned trust
root. Labels contain $1$--$255$ printable non-space ASCII bytes and contexts
have at most $65536$ bytes. Thus
\[
 |I|=107+|a|+|c|\in[108,65898],\qquad |I\Vert\mathtt{01}|\le65899.
\]
Length parsing recovers every effective field. The descriptor is
$(\Gamma(r_0),I)$ across all owners of the same protected contribution.
Runtime identities and owner identifiers retain their separate lifecycle
roles. The effective digest bytes determine their respective info fields.

\paragraph{Owner correspondence.}
Fix a protected contribution family with common value $S$ and preserve the
actual metadata of each source and derived owner. Represent a derived owner
by $(\mathit{runtime},\mathit{live},I)$ and a reference to the common key
entry $K_I$. Every accepted derivation creates a new owner; equal $I$ values
reference the same entry. Closing one derived owner changes its own live
flag. Closing a source owner changes its future derivation access; existing
derived owners retain their stored key. Runtime closure makes later open checks fail. An admitted derivation or an
export whose open check already succeeded follows its actual completion
path. Previously exported copies remain available to their recipients.

These transitions follow directly from the source's separate
\texttt{Option} fields, shared runtime state and admission checks. The inspected FFI derivation path additionally retains reservation,
publication and parent-revocation metadata. A first private KDF computation allocates the descriptor
entry even if its later handle publication fails; a retry retains that key
entry while obtaining its actual new owner. The Rust in-flight bound and the
FFI live-handle bound are preserved as metadata. The proof's lifetime bound
$N$ counts distinct complete infos.

\begin{proposition}[Native SDK key-table correspondence]
\label{prop:sdk-purpose-table}
For one initially independent uniform protected contribution, let the
family's source owners use the Rust derivation interface above, together with
derived-owner export, close and runtime lifecycle operations. The parent
contribution and internal Extract value remain within their stated KDF
interfaces. Let $N$ bound distinct infos before execution and $L$ bound their
Expand-message length. Then the entire owned-API experiment differs from
one using independent uniform keys $K_I$, with all aliases and owner states
preserved, by output-bit gap at most
\begin{equation}
\begin{split}
 \Delta_{\rm sdk}=\min\{1,\;&
 \varepsilon_{\rm Hjoint}(1,0,B+2+(2+\ell(L))N,32)\\
 &+\varepsilon_{\rm Hrow}(N,B+2,L)\}.
\end{split}
 \label{eq:sdk-purpose-error}
\end{equation}
The budgets cover complete preparation, other real computations and the
terminal test. At full SDK input limits, $L=65899$ and $\ell(L)=1031$.
\end{proposition}
\begin{proof}
Normalize the fixed salt as above. The first-block HKDF definition~\cite{rfc5869}
gives $P=F(r_0,S)$ followed by $F(P,I\Vert\mathtt{01})$; the source keeps
its PRK internal.
An induction over admitted calls preserves each owner state and every
public status: validation and admission use available metadata; derivation
creates a fresh owner referring to its descriptor's deterministic value;
exports apply that owner's current guards; close and publication operate
on their existing owner identities. Repeated key values can therefore be
memoized while all native owner transitions continue to execute.
Apply Lemma~\ref{lem:entry-information-length} with $D=C=1$.
The class/row comparison changes only these internal key values. All
field-dependent choices, exports and lifecycle responses are recomputed
in their own experiment by the same controller.
\end{proof}

For the existing five-purpose source workload, the label has 13 bytes and
its context 24, giving $|I|=144$, $L=145$ and $\ell(L)=4$. Its two peer
owners use five complete descriptors when their effective policy bindings
agree. Thus $N=5$ gives column ordinary budget $B+32$ and row budget $B+2$.
The table comparison permits disclosure of all derived entries and
establishes their joint pseudorandom separation. In particular, the
probability of a key collision among exports from distinct allocated infos,
with at most $N$ total infos, is at most $\min\{1,\binom N2 2^{-256}+\Delta_{\rm sdk}\}$.

\paragraph{A selected key for one application encryption.}
An application may mark a new complete descriptor before its first private
KDF evaluation, retain its key privately, and use it in the one-key,
one-encryption AEAD experiment defined in
Section~\ref{sec:joint-continuation}, with that experiment's nonce,
equal-length messages, associated data and decryption restriction. Every source owner and derived alias of that descriptor routes
to the same key-use interface, with its own native liveness checks. Other
derived entries may be exported. Let $Z$ be the qualified signed score,
with qualification computed from the available classical view and ledger
and value zero when the marked challenge is absent. Fresh table allocation
then supplies the uniform AEAD key, giving
\[
 |\mathbb E Z|\le\min\{1,\,2\Delta_{\rm sdk}
                                  +2\varepsilon_{\rm AEAD}\}.
\]
Fix each application challenge bit for the KDF comparison and average the
two score bounds. In the final AEAD game, use the adversary's guess on
qualified runs and an independent fair guess otherwise. A preceding
complete-tuple ContextBound comparison contributes its established
$2\Phi_{256}(d,Q_H)$ term when the stopped execution realizes that source
family and available controller.

This consequence concerns the exported purpose key's declared application
use. The durable rekey driver retains its own 97-byte pending-root
information and established source invariant; the SDK table provides a
separate concrete application of the finite entry proof.
